\documentclass[11pt]{article}
\usepackage[a4paper,margin=23mm]{geometry}
\usepackage{amsmath,amssymb,amsthm,hyperref}
\newtheorem{theorem}{Theorem}
\newtheorem{corollary}[theorem]{Corollary}
\newtheorem{lemma}[theorem]{Lemma}
\newtheorem{remark}[theorem]{Remark}
\title{GST Optimal-Dimensional Whole-Space Incidence Realization and Exact Finite-Trace Hodge Residual}
\date{23 October 2026}
\begin{document}\maketitle
\begin{abstract}
Given an arbitrary original smooth projective complex d-fold X,
we determine the SHARP minimal possible complex dimension
of any Stein affine-space torsor over X: its affine-fiber rank
must be at least d. The bound is attained for every X by
pulling back the genuine incidence-complement affine torsor
over P^d along a finite Noether-normalization map X->P^d.
This provides a UNIVERSAL realized GST whole-space
fold of minimal total complex dimension 2d, with exact
ordinary/compact-support Hodge signatures and the actual
rank-d higher-causal translation geometry.
Next we construct in rational Betti AND true Chow
the exact finite-cover transfer projector onto the
algebraic ambient Hodge direction. We show that every
prescribed rational (p,p) alpha splits canonically for
each finite projection into a constructed genuine algebraic
cycle and a trace-zero Hodge residual. We prove that
pure incidence/fiber-rank manipulations cannot manufacture
the latter: the finite-transfer original-X residual
is invariant under all base-induced operations.
The work is actual all-X mathematics; no special
example is extrapolated to the universal Hodge target.
\end{abstract}

\section{Exact geometry and original-X source target}
Let X be any connected smooth projective complex
variety, $\dim_\mathbb C X=d\ge1$.
The target of the full GST program remains:
for every p with $0\le p\le d$ and every rational
Hodge class $\alpha\in H^{2p}(X,\mathbb Q)\cap H^{p,p}(X)$,
construct a GENUINE
$Z_\alpha\in CH^p(X)_\mathbb Q$
with actual canonical rational Betti
$\mathrm{cl}_X(Z_\alpha)=\alpha$.
No target Chow class is ever an input below.

A genuine affine-space torsor of rank r
means a Zariski-locally trivial morphism
$q:J\to X$ with fibers $\mathbb A^r_\mathbb C$,
under some actual algebraic vector bundle
of rank r. Its topology is a genuine
affine-fiber fold, not a formal
higher-cell dimension declaration.

\begin{theorem}[UNIVERSAL SHARP RANK LOWER BOUND]
If the analytification of the total space J
of a rank-r affine-space torsor
$q:J\to X$ is Stein, then
\[
 \boxed{r\ge d,\qquad
        \dim_\mathbb C J=d+r\ge2d.}                  \tag{1}
\]
This holds on EVERY smooth projective
complex d-fold X without any
Hodge-algebraicity hypothesis.
\end{theorem}
\begin{proof}
Each affine-space fiber is contractible;
because q is locally trivial analytically,
it is a homotopy equivalence. In particular,
\[
 H^{2d}(J,\mathbb Q)\cong H^{2d}(X,\mathbb Q)
              \cong\mathbb Q
\]
by the nonzero compact complex orientation
fundamental class of connected X.
A Stein manifold of complex dimension n
has the homotopy type of a CW complex
of REAL dimension at most n,
by the Andreotti--Frankel/Morse
handle-index theorem for Stein manifolds.
Applied to J of complex dimension d+r,
it yields $H^k(J,\mathbb Q)=0$ for
$k>d+r$. Since $H^{2d}(J,\mathbb Q)$
is nonzero, $2d\le d+r$, proving (1).
Neither rational Hodge type nor the
Hodge conjecture is used here.
\end{proof}

\begin{theorem}[ACTUAL SHARP MINIMAL-RANK FOLD: EVERY original X]
Every smooth projective complex d-fold X
admits a smooth AFFINE, Stein-analytic
rank-d affine-space torsor
\[
 \boxed{
 q_d:J_d^{\min}(X)\to X,\qquad
 \dim_\mathbb C J_d^{\min}(X)=2d.
 }                                                       \tag{2}
\]
The total-space algebraic scheme has
the explicit genuine incidence equation
\[
 \boxed{
 J_d^{\min}(X)=
 \{(x,H)\in X\times(\mathbb P^d)^\vee:
                    f(x)\notin H\},
 }                                                       \tag{3}
\]
where $f:X\to\mathbb P^d$ is ANY finite
surjective projective Noether-normalization map.
It is the affine torsor under the
rank-d bundle
\[
 \boxed{
 V_f=f^*\operatorname{Hom}
       (Q_{\mathbb P^d},O_{\mathbb P^d}(-1)).
 }                                                       \tag{4}
\]
For all natural k, replacing this space
by $J_d^{\min}(X)\times\mathbb A^k$
produces a genuine Stein affine
rank-(d+k) fold; none with rank<d
can be Stein.
\end{theorem}
\begin{proof}
For any positive-dimensional projective
d-dimensional X, homogeneous
Noether normalization for its
projective homogeneous coordinate
ring, or generic finite linear
projection, supplies a finite
surjective morphism
$f:X\to\mathbb P^d$ of degree
$\delta>0$. The standard incidence
complement
\[
 U_d=\{(y,H)\in\mathbb P^d\times
 (\mathbb P^d)^\vee:y\notin H\}
\]
is the complement of a
$(1,1)$-divisor in the product.
Since $O(1,1)$ is ample,
this complement is an AFFINE
algebraic variety.
The natural projection $U_d\to\mathbb P^d$
has fiber $\mathbb P^d\setminus H$
isomorphic to $\mathbb A^d$;
as the space of retractions of
$O(-1)\subset O^{d+1}$ it is
a Zariski-locally trivial torsor
under $V_{\mathbb P^d}$ in (4).
Base change along f gives exactly
$J_d^{\min}(X)=X\times_{\mathbb P^d}U_d$.
The projection $J_d^{\min}(X)\to U_d$
is finite (finite morphisms are
stable under base change).
A finite morphism over an affine
scheme has AFFINE source, so
$J_d^{\min}(X)$ is affine.
It is smooth as an affine-space
torsor over smooth X and has
complex dimension 2d.
Its analytification is Stein.
The torsor bundle pulls back
to (4).
Products with any affine A^k
are affine and remain actual
torsors; the preceding theorem
proves minimality.
\end{proof}

\begin{corollary}[Universal whole-Hodge causal fold law at the SHARP minimum]
Every original X has its own
geometrically optimal entire-space
realization q_d in (2), whose ACTUAL
rational mixed Hodge structure and
compact-support Hodge polynomial obey
\[
 \boxed{
 q_d^*:H^j(X,\mathbb Q)
    \overset{\sim}{\longrightarrow}
      H^j(J_d^{\min}(X),\mathbb Q)
 \ \text{as rational MHS, for all }j,}
                                                               \tag{5}
\]
\[
 \boxed{
 E_c(J_d^{\min}(X);u,v)
    =(uv)^d\,E_c(X;u,v).
 }                                                       \tag{6}
\]
The first identifies the entire
ordinary Hodge structure in EVERY
grade; the second shifts compact
weights by exactly d Tate degrees.
The actual torsor's higher-causal
vertical tangent has rank d,
and its relative curvature tensor
lies in
$\Omega_{J/X}^2\otimes\mathrm{End}(H)$.
For d>=2 this sheaf is nonzero
as a vector-bundle sheaf;
for d=1 it vanishes identically.
Thus the minimal REAL spatial
complex dimension cannot force
nonabelian curvature to vanish
by a rank-one degree argument
when d>=2.
\end{corollary}
\begin{proof}
Algebraic affine-space bundles
are A1-homotopy invariant on
rational singular cohomology
with its natural mixed Hodge
structure (equivalently, use the
algebraic Leray sheaf R^j q_*Q,
which vanishes above zero
because each geometric affine
fiber is contractible).
The compact-support
Hodge--Deligne polynomial is
multiplicative under a
Zariski locally trivial
algebraic A^d fibration,
whose fiber has $E_c=(uv)^d$.
The relative cotangent bundle
is rank d, so its exterior
square has rank d(d-1)/2;
this vanishes exactly for
d<=1 and is nonzero for d>=2.
\end{proof}

\section{A TRUE finite-cover Chow return and exact trace-zero Hodge residual}
Write $f:X\to\mathbb P^d$ for
the actual finite morphism
used in the minimal incidence
space, $\delta=\deg f$,
$h=c_1(f^*O_{\mathbb P^d}(1))$.
Let $H=c_1^{CH}(f^*O(1))$ denote
the true original-X
hyperplane Chow divisor.
The map f is finite and flat:
the target P^d is smooth,
the source X is Cohen--Macaulay,
the fibers are zero-dimensional,
and the miracle-flatness
criterion applies.

\begin{theorem}[Actual rational Chow and Betti finite-trace PROJECTOR]
For EVERY p, 0<=p<=d,
there is an actual algebraic
correspondence-type Chow operator
\[
 \boxed{
 T^{CH}_{f,p}=
 \frac1\delta f^*f_*:
 CH^p(X)_\mathbb Q\to CH^p(X)_\mathbb Q
 }                                                       \tag{7}
\]
using the GENUINE flat pullback
and finite proper pushforward.
It is idempotent, with image
$\mathbb Q H^p$ modulo any
Chow relations; since $\int_XH^d
=\delta$, $H^p$ is nonzero
in the relevant rational Chow.
The corresponding Betti operator
\[
 \boxed{
 T_{f,p}=\frac1\delta f^*f_*:
 H^{2p}(X,\mathbb Q)\to H^{2p}(X,\mathbb Q)
 }                                                       \tag{8}
\]
is idempotent, a morphism of
rational Hodge structures
of type (0,0) in degree 2p,
and has exact image $\mathbb Q h^p$.
For EVERY rational Hodge alpha
define the ACTUAL rational
number
\[
 c_\alpha=\int_X\alpha\,h^{d-p}\in\mathbb Q.
\]
Then the identity of ORIGINAL-X
rational Betti classes is
\[
 \boxed{
 \alpha=\underbrace{\frac{c_\alpha}
  {\delta}h^p}_{\mathrm{cl}_X^{can}
  ((c_\alpha/\delta)H^p)}
 +\underbrace{\left(
   \alpha-\frac{c_\alpha}{\delta}h^p
             \right)}_{\alpha_f^0},
 \qquad
 f_*(\alpha_f^0)=0 .
 }                                                       \tag{9}
\]
The first term has the displayed
ACTUAL independently constructed
original-X Chow source. The
second is rational of type (p,p)
and lives in the TRUE trace-zero
subspace of the whole spatial
finite-fold geometry.
\end{theorem}
\begin{proof}
The finite flat projection
formula says
$f_*f^*\eta=\delta\eta$
for cohomology and Chow
with rational coefficients.
Therefore
$(f^*f_*/\delta)^2
=f^*(f_*f^*)f_*/\delta^2
=f^*f_*/\delta$.
The cohomology and Chow rings
of P^d contain only the
one-dimensional group
in codimension p generated by
$\zeta^p$,
where $\zeta=c_1(O(1))$.
Hence both projector images
are spanned by $f^*\zeta^p=h^p$
(or H^p in Chow).
To determine the coefficient
of $f_*\alpha\in H^{2p}(P^d)$,
integrate against $\zeta^{d-p}$:
\[
 \int_{\mathbb P^d}
 f_*\alpha\,\zeta^{d-p}
 =\int_X\alpha\,f^*\zeta^{d-p}
 =c_\alpha .
\]
Thus $f_*\alpha=c_\alpha\zeta^p$
and (9) follows. The canonical
cycle-class map commutes with
actual finite flat pullback
and proper pushforward.
Because f is algebraic, f^* and
f_* respect rational Hodge
bidegrees after the proper
degree shifts (none here, since
finite). Hence the residual is
Hodge whenever alpha is Hodge.
\end{proof}

\begin{corollary}[Exact original-X residual source equivalence]
For EVERY original X and every
finite projective f as above,
\[
 \boxed{
 \alpha\in\mathrm{im}(\mathrm{cl}^{can}_X)
 \quad\Longleftrightarrow\quad
 \alpha_f^0\in
   \mathrm{im}(\mathrm{cl}^{can}_X).
 }                                                       \tag{10}
\]
In particular for p=0 and p=d,
the trace-zero rational Hodge
residual vanishes for every alpha,
and the original-X canonical
algebraic source is completely
constructed by the first term
of (9). In middle grades the
trace-zero subspace is the exact
new geometry that any genuine
universal source birth must
represent, not a change of
notation for an ordinary
hyperplane class.
\end{corollary}
\begin{proof}
The first term in (9) is an
ACTUAL rational Chow cycle on
original X for every input alpha,
without alpha being algebraic.
Thus alpha is algebraic iff its
difference with that term is
algebraic. In degrees zero and
2d the original rational
cohomology groups are
one-dimensional for connected
smooth projective X, and are
already spanned by h^p, so the
trace-zero kernel vanishes.
\end{proof}

\section{Independent source creation inside the trace-zero geometry}
Let $K_X$ denote the actual
algebraic canonical divisor
and let $R_f$ denote the genuine
effective ramification divisor
of the finite flat generically
separable map f.
The Riemann--Hurwitz canonical
divisor identity gives
\[
 \boxed{
 R_f=K_X+(d+1)H\quad
          \text{in }CH^1(X)_\mathbb Q .
 }                                                       \tag{11}
\]
In ANY positive codimension p
construct the ACTUAL Chow source
\[
 \boxed{
 B^{CH}_{f,p}=
 R_f\cdot H^{p-1}
 -\frac{\int_X c_1(R_f)h^{d-1}}
        {\delta}\,H^p
 \in CH^p(X)_\mathbb Q .
 }                                                       \tag{12}
\]
This formula is always defined
without a chosen Hodge target
cycle. Its canonical rational
Betti class lies in
$\ker f_*$ by the exact
projector theorem. It can be
NONZERO when the class of
$R_f$ has a nonambient component
that survives multiplying by
h^(p-1). It can vanish,
notably when Pic(X)_Q is
generated by H, so it does
NOT give universal spanning
of all trace-zero Hodge
directions. This is a genuinely
constructed nonambient source
from the original finite cover
and its branch geometry, distinct
from the hyperplane-only factory.

\section{Actual all-N higher-causal constraint and the remaining universal mechanism}
The minimal whole-world GST fold
exists for EVERY original smooth
projective X, and the first
rational Hodge-preserving
spatial source projection is
solved by (9). Its actual
higher-causal vector-bundle
groupoid has vertical rank d.
The fold's own topological
trace-zero sector can carry
rational Hodge classes that
are not exhausted by the
finite branch source (12).
The native K0 algebraization
quotient described in the
existing manuscript
\texttt{GST\_stein\_affine\_line\_causal\_flatness\_universal\_chow\_2026\_10\_23.tex}
is preserved under ANY
affine-space torsor:
\[
 \frac{K_0^{hol}(J)_\mathbb Q}
 {\operatorname{im}K_0^{alg}(J)_\mathbb Q}
 \cong\bigoplus_p
 \frac{H^{2p}(X,\mathbb Q)}
 {\operatorname{im}\mathrm{cl}_X^{can,p}} .
 \tag{13}
\]
Equation (13) is exact for the
OPTIMAL J in (2), just as
for larger unnecessary ambient
ranks. Thus the geometric
problem is not resolved by
adding artificial extra affine
dimensions, nor by taking
finite traces back to P^d:
(9) shows precisely where
the original-X Hodge residual
persists. A genuinely new
higher-causal construction
would have to create an actual
algebraic lift for that
trace-zero rational Hodge
class, independently of
assuming it in advance.

The achieved theorem is the
SHARP universal spatial
dimension law, its geometric
realization, the true
finite-cover Chow projector,
and the original-X trace-zero
source decomposition.
This is verified for all X by
topology, affine geometry,
finite morphism geometry
and actual rational Chow
push-pull. It does not replace
the residual source with a
conditional flatness hypothesis
and claim a full solution.
\end{document}